% !TeX root = part2.tex

\documentclass[../final.tex]{subfiles}

\begin{document}

% ============================================================
% Практическое занятие 2
% ============================================================

\practice{2}{10.09.2026}{Взаимодействие с веществом}


% ============================================================
% Основные механизмы взаимодействия с веществом
% ============================================================

\rtheory{Основные механизмы взаимодействия с веществом}

При прохождении частиц через вещество характер взаимодействия
существенно зависит от того, является частица заряженной
или нейтральной.


% ------------------------------------------------------------
\subsection*{Заряженные частицы}
% ------------------------------------------------------------

К заряженным частицам относятся, например,

\begin{equation*}
    e^\pm,
    \qquad
    p,
    \qquad
    \mu^\pm,
    \qquad
    K^\pm,
    \qquad
    \pi^\pm.
\end{equation*}

При прохождении через вещество для них существенны следующие процессы:

\begin{itemize}
    \item ионизационные потери;
    \item тормозные потери;
    \item многократное рассеяние;
    \item черенковское излучение;
    \item переходное излучение.
\end{itemize}

Ионизационные и тормозные процессы приводят прежде всего
к \rresult{потере энергии} частицы.

Многократное рассеяние приводит главным образом к
\rresult{изменению траектории} движения.


% ------------------------------------------------------------
\subsection*{Нейтральные частицы}
% ------------------------------------------------------------

К рассматриваемым нейтральным частицам относятся

\begin{equation*}
    n,
    \qquad
    \gamma.
\end{equation*}

Для $\gamma$-квантов основными процессами взаимодействия
с веществом являются

\begin{itemize}
    \item фотоэффект;
    \item эффект Комптона;
    \item рождение электрон-позитронных пар.
\end{itemize}


% ============================================================
% Классическая оценка передачи энергии электрону
% ============================================================

\rtheory{Классическая оценка передачи энергии электрону}


% ------------------------------------------------------------
\task[1]{Пролёт тяжёлой заряженной частицы мимо электрона}

\condition
% ------------------------------------------------------------

Рассмотрим тяжёлую частицу с зарядом $Ze$,
движущуюся с постоянной скоростью $v$ мимо атомного электрона.

Пусть минимальное расстояние от траектории частицы
до электрона равно $b$.

Требуется определить энергию

\begin{equation*}
    \Delta E,
\end{equation*}

переданную электрону при таком взаимодействии.


\begin{rbox}[left=18pt,right=18pt,top=14pt,bottom=10pt]{[]}
    \centering

    \begin{tikzpicture}[
        >=Latex,
        every node/.style={text=rtext}
    ]

        % -----------------------------------------------------
        % Оси
        % -----------------------------------------------------

        \draw[
            ->,
            draw=rtext,
            line width=0.8pt
        ]
        (-4.3,0) -- (4.4,0)
        node[right] {$x$};

        \draw[
            ->,
            draw=rtext,
            line width=0.8pt
        ]
        (0,-0.5) -- (0,3.4)
        node[above] {$y$};


        % -----------------------------------------------------
        % Электрон
        % -----------------------------------------------------

        \coordinate (E) at (0,2.2);

        \fill[rc3] (E) circle (2.5pt);

        \node[
            above right=2pt
        ] at (E) {$e^-$};


        % -----------------------------------------------------
        % Налетающая частица
        % -----------------------------------------------------

        \coordinate (P) at (-2.7,0);

        \fill[rc3] (P) circle (2.5pt);

        \node[
            below=6pt
        ] at (P) {$Ze$};


        % -----------------------------------------------------
        % Вектор скорости прямо на оси x
        % -----------------------------------------------------

        \draw[
            ->,
            draw=rc3,
            line width=1.25pt
        ]
        (-2.62,0) -- (-1.05,0);

        \node[
            above=5pt
        ] at (-1.78,0) {$\vec v$};


        % -----------------------------------------------------
        % Прицельный параметр b
        % -----------------------------------------------------

        \draw[
            densely dashed,
            draw=rtextmuted,
            line width=0.7pt
        ]
        (0,0) -- (E);

        \node[
            right=5pt
        ] at (0,1.05) {$b$};


        % -----------------------------------------------------
        % Линия между частицей и электроном
        % -----------------------------------------------------

        \draw[
            draw=rtext,
            line width=0.85pt
        ]
        (P) -- (E);


        % -----------------------------------------------------
        % Сила
        % -----------------------------------------------------

        \draw[
            ->,
            draw=rc3,
            line width=1.15pt
        ]
        (E) -- ++(-1.45,-1.18);

        \node[
            above left=1pt
        ] at (-0.85,1.55) {$\vec F$};


        % -----------------------------------------------------
        % Угол theta
        % -----------------------------------------------------

        \draw[
            draw=rtext,
            line width=0.8pt
        ]
        ($(E)+(0,-0.62)$)
        arc[
            start angle=-90,
            end angle=-140,
            radius=0.62
        ];

        \node at (-0.31,1.48) {$\theta$};

    \end{tikzpicture}

    \captionsetup{
        font=footnotesize,
        labelfont=bf,
        skip=5pt
    }

    \captionof{figure}{
        Пролёт тяжёлой заряженной частицы с зарядом $Ze$
        мимо атомного электрона.
        Параметр $b$ определяет минимальное расстояние
        от траектории частицы до электрона.
    }
    \label{fig:matter_collision_geometry}
\end{rbox}


Из геометрии рисунка расстояние между налетающей частицей
и электроном равно

\begin{equation*}
    r
    =
    \frac{b}{\cos\theta}.
\end{equation*}

Кулоновская сила взаимодействия имеет модуль

\begin{equation*}
    F
    =
    \frac{Ze^2}{r^2}
    =
    \frac{Ze^2}{
        \left(
            b/\cos\theta
        \right)^2
    }.
\end{equation*}

Нас интересует поперечная составляющая силы:

\begin{equation*}
    F_y
    =
    F\cos\theta
    =
    \frac{Ze^2}{
        \left(
            b/\cos\theta
        \right)^2
    }
    \cos\theta
    =
    \frac{Ze^2}{b^2}
    \cos^3\theta.
\end{equation*}

При равномерном движении налетающей частицы

\begin{equation*}
    x=vt.
\end{equation*}

Из геометрии

\begin{equation*}
    \tan\theta
    =
    \frac{vt}{b}.
\end{equation*}

Следовательно,

\begin{equation*}
    \cos\theta
    =
    \frac{1}{
        \sqrt{
            1+\tan^2\theta
        }
    }
    =
    \frac{1}{
        \sqrt{
            1+
            \left(
                \dfrac{vt}{b}
            \right)^2
        }
    }.
\end{equation*}

Изменение поперечной компоненты импульса электрона определяется
соотношением

\begin{equation*}
    dp_y
    =
    F_y\,dt.
\end{equation*}

Поэтому полный поперечный импульс, переданный электрону,

\begin{equation*}
    p_y
    =
    \int_{-\infty}^{+\infty}
    F_y\,dt.
\end{equation*}

Подставим выражение для силы:

\begin{equation*}
    p_y
    =
    \int_{-\infty}^{+\infty}
    \frac{Ze^2}{b^2}
    \cos^3\theta\,dt
    =
    \frac{Ze^2}{b^2}
    \int_{-\infty}^{+\infty}
    \left[
        1+
        \left(
            \frac{vt}{b}
        \right)^2
    \right]^{-3/2}
    dt.
\end{equation*}

Преобразуем выражение в скобках:

\begin{equation*}
    1+
    \frac{v^2t^2}{b^2}
    =
    \frac{v^2}{b^2}
    \left[
        \left(
            \frac{b}{v}
        \right)^2
        +
        t^2
    \right].
\end{equation*}

Поэтому

\begin{equation*}
    \left(
        1+
        \frac{v^2t^2}{b^2}
    \right)^{-3/2}
    =
    \frac{b^3}{v^3}
    \left[
        \left(
            \frac{b}{v}
        \right)^2
        +
        t^2
    \right]^{-3/2}.
\end{equation*}

Тогда

\begin{equation*}
    p_y
    =
    \frac{Ze^2b}{v^3}
    \int_{-\infty}^{+\infty}
    \frac{dt}{
        \left[
            \left(
                \dfrac{b}{v}
            \right)^2
            +
            t^2
        \right]^{3/2}
    }.
\end{equation*}

Используем интеграл

\begin{equation*}
    \int
    \frac{dt}{
        \left(
            a^2+t^2
        \right)^{3/2}
    }
    =
    \frac{t}{
        a^2
        \sqrt{
            a^2+t^2
        }
    }.
\end{equation*}

При

\begin{equation*}
    a
    =
    \frac{b}{v}
\end{equation*}

получаем

\begin{equation*}
    p_y
    =
    \frac{Ze^2b}{v^3}
    \left.
    \frac{t}{
        \left(
            \dfrac{b}{v}
        \right)^2
        \sqrt{
            \left(
                \dfrac{b}{v}
            \right)^2+t^2
        }
    }
    \right|_{-\infty}^{+\infty}.
\end{equation*}

При $t\rightarrow\pm\infty$

\begin{equation*}
    \frac{t}{
        \sqrt{
            a^2+t^2
        }
    }
    \longrightarrow
    \pm1.
\end{equation*}

Поэтому

\begin{equation*}
    p_y
    =
    \frac{Ze^2b}{v^3}
    \frac{v^2}{b^2}
    \left[
        1-(-1)
    \right]
    =
    \rboxed{
        \frac{2Ze^2}{bv}
    }.
\end{equation*}

Полученная электроном кинетическая энергия равна

\begin{equation*}
    \Delta E
    =
    \frac{p_y^2}{2m_e}.
\end{equation*}

Подставляя найденный импульс,

\begin{equation*}
    \Delta E
    =
    \frac{1}{2m_e}
    \left(
        \frac{2Ze^2}{bv}
    \right)^2
    =
    \frac{4Z^2e^4}{
        2m_eb^2v^2
    }
    =
    \frac{2Z^2e^4}{
        m_eb^2v^2
    }.
\end{equation*}

Так как

\begin{equation*}
    v=\beta c,
\end{equation*}

окончательно

\begin{equation*}
    \rboxed{
        \Delta E
        =
        \frac{
            2Z^2e^4
        }{
            b^2\beta^2m_ec^2
        }
    }.
\end{equation*}

Для характерного атомного расстояния

\begin{equation*}
    b
    =
    10^{-8}\,\text{см}
\end{equation*}

получается оценка

\begin{equation*}
    \rboxed{
        \Delta E
        \simeq
        8\cdot10^{-4}
        \left(
            \frac{Z}{\beta}
        \right)^2
        \text{эВ}
    }.
\end{equation*}


\remark[Домашнее задание.]
Оценить потери энергии для

\begin{equation*}
    \alpha,
    \qquad
    p,
    \qquad
    \pi
\end{equation*}

при кинетической энергии

\begin{equation*}
    T
    =
    1\,\text{МэВ}.
\end{equation*}


% ============================================================
% Задачи на ионизационные потери
% ============================================================

\rtheory{Ионизационные потери}

Для дальнейших оценок используем формулу ионизационных потерь
тяжёлой заряженной частицы

\begin{equation*}
    -
    \frac{dE}{dx}
    =
    4\pi
    N_A
    r_e^2
    m_ec^2
    z^2
    \frac{Z}{A}
    \frac{1}{\beta^2}
    \left[
        \ln
        \left(
            \frac{
                2m_ec^2
                \gamma^2\beta^2
            }{
                I
            }
        \right)
        -
        \beta^2
        -
        \frac{\delta}{2}
    \right].
\end{equation*}

Здесь

\begin{equation*}
    \beta
    =
    \frac{v}{c},
    \qquad
    \gamma
    =
    \frac{1}{
        \sqrt{
            1-\beta^2
        }
    }.
\end{equation*}

Для численных оценок используется приближённая форма

\begin{equation*}
    \rboxed{
        -
        \frac{dE}{dx}
        \simeq
        0.3
        \frac{
            Zz^2
        }{
            A\beta^2
        }
        \left[
            11.2
            +
            \ln
            \left(
                \frac{
                    (\beta\gamma)^2
                }{
                    Z
                }
            \right)
            -
            \beta^2
        \right]
    }.
\end{equation*}

Характерное минимальное значение удельных ионизационных потерь:

\begin{equation*}
    \rboxed{
        -
        \left(
            \frac{dE}{dx}
        \right)_{\min}
        \simeq
        2\,
        \frac{
            \text{МэВ}
        }{
            \text{г}/\text{см}^2
        }
    }.
\end{equation*}


% ------------------------------------------------------------
\task[2]{Потери энергии в углероде и свинце}

\condition
% ------------------------------------------------------------

Сравним ионизационные потери одной и той же частицы
при прохождении через углерод и свинец.

Для углерода

\begin{equation*}
    Z_{\mathrm C}=6,
    \qquad
    A_{\mathrm C}\simeq12,
\end{equation*}

поэтому

\begin{equation*}
    \frac{Z_{\mathrm C}}{A_{\mathrm C}}
    \simeq
    \frac{6}{12}
    =
    \frac{1}{2}.
\end{equation*}

Для свинца

\begin{equation*}
    Z_{\mathrm{Pb}}=82,
    \qquad
    \frac{Z_{\mathrm{Pb}}}{A_{\mathrm{Pb}}}
    \simeq
    \frac{2}{5}.
\end{equation*}

При одинаковых $z$ и $\beta$ получаем

\begin{equation*}
    \frac{
        \left(
            -\dfrac{dE}{dx}
        \right)_{\mathrm C}
    }{
        \left(
            -\dfrac{dE}{dx}
        \right)_{\mathrm{Pb}}
    }
    =
    \frac{
        \dfrac{1}{2}
        \left[
            11.2
            +
            \ln
            \left(
                \dfrac{
                    (\beta\gamma)^2
                }{
                    6
                }
            \right)
            -
            \beta^2
        \right]
    }{
        \dfrac{2}{5}
        \left[
            11.2
            +
            \ln
            \left(
                \dfrac{
                    (\beta\gamma)^2
                }{
                    82
                }
            \right)
            -
            \beta^2
        \right]
    }.
\end{equation*}

Рассмотрим протон с кинетической энергией

\begin{equation*}
    T
    =
    10\,\text{МэВ}.
\end{equation*}

Его полная энергия

\begin{equation*}
    E
    =
    E_0+T
    =
    \gamma m_pc^2.
\end{equation*}

При

\begin{equation*}
    m_pc^2
    \simeq
    938\,\text{МэВ}
\end{equation*}

имеем

\begin{equation*}
    E
    \simeq
    948\,\text{МэВ}.
\end{equation*}

Следовательно,

\begin{equation*}
    \gamma
    =
    \frac{948}{938}.
\end{equation*}

Используем

\begin{equation*}
    \gamma
    =
    \frac{1}{
        \sqrt{
            1-\beta^2
        }
    }.
\end{equation*}

Отсюда

\begin{equation*}
    \beta^2
    =
    1-
    \frac{1}{\gamma^2}
    \simeq
    0.021.
\end{equation*}

После подстановки получаем

\begin{equation*}
    \rboxed{
        \frac{
            \left(
                -\dfrac{dE}{dx}
            \right)_{\mathrm C}
        }{
            \left(
                -\dfrac{dE}{dx}
            \right)_{\mathrm{Pb}}
        }
        \simeq
        2.3
    }.
\end{equation*}


% ------------------------------------------------------------
\task[3]{Зависимость потерь энергии от глубины}

\condition
% ------------------------------------------------------------

Рассмотрим протон с начальной кинетической энергией

\begin{equation*}
    T
    \simeq
    1\,\text{ГэВ},
\end{equation*}

влетающий в вещество.

Необходимо определить характер зависимости

\begin{equation*}
    \frac{dE}{d\ell}
    =
    \frac{dE}{d\ell}(\ell).
\end{equation*}


\begin{rbox}[left=18pt,right=18pt,top=14pt,bottom=10pt]{[]}
    \centering

    \begin{tikzpicture}[
        >=Latex,
        every node/.style={text=rtext}
    ]

        % -----------------------------------------------------
        % Границы вещества
        % -----------------------------------------------------

        \draw[
            draw=rtextmuted,
            line width=0.8pt
        ]
        (-0.4,-1.25) -- (-0.4,1.25);

        \draw[
            draw=rtextmuted,
            line width=0.8pt
        ]
        (3.6,-1.25) -- (3.6,1.25);


        % -----------------------------------------------------
        % Протон до входа
        % -----------------------------------------------------

        \fill[rc3] (-3.3,0) circle (2.5pt);

        \node[
            above=5pt
        ] at (-3.3,0) {$p$};

        \draw[
            ->,
            draw=rc3,
            line width=1.2pt
        ]
        (-3.2,0) -- (-0.45,0);

        \node[
            below=6pt
        ] at (-1.8,0) {$T\simeq1\,\text{ГэВ}$};


        % -----------------------------------------------------
        % Траектория внутри вещества
        % -----------------------------------------------------

        \draw[
            ->,
            draw=rtext,
            line width=0.9pt
        ]
        (-0.4,0) -- (2.9,0);

        \fill[rc3] (2.95,0) circle (2.3pt);

        \node[
            above
        ] at (1.3,0) {$\ell$};

        \node[
            below
        ] at (1.6,-0.35) {вещество};

    \end{tikzpicture}

    \captionsetup{
        font=footnotesize,
        labelfont=bf,
        skip=5pt
    }

    \captionof{figure}{
        Прохождение протона с начальной кинетической энергией
        порядка $1\,\text{ГэВ}$ через вещество.
    }
    \label{fig:proton_matter_depth}
\end{rbox}


Для оценки примем

\begin{equation*}
    \gamma
    \simeq
    2.
\end{equation*}

Из связи между $\gamma$ и $\beta$

\begin{equation*}
    \gamma
    =
    \frac{1}{
        \sqrt{
            1-\beta^2
        }
    }
\end{equation*}

следует

\begin{equation*}
    \beta
    =
    \sqrt{
        1-
        \frac{1}{\gamma^2}
    }
    =
    \sqrt{
        1-
        \frac14
    }
    =
    \frac{\sqrt3}{2}.
\end{equation*}

Следовательно,

\begin{equation*}
    \beta^2
    =
    \frac34.
\end{equation*}

Запишем приближённое выражение для ионизационных потерь:

\begin{equation*}
    -
    \frac{dE}{d\ell}
    \simeq
    0.3
    \frac{
        Zz^2
    }{
        A\beta^2
    }
    \left[
        11.2
        +
        \ln
        \left(
            \frac{
                (\beta\gamma)^2
            }{
                Z
            }
        \right)
        -
        \beta^2
    \right].
\end{equation*}

Аргумент логарифма в рассматриваемой области меняется
примерно в пределах от $2$ до $3.5$, поэтому логарифмический
член можно оценивать как

\begin{equation*}
    \ln(2\text{--}3.5).
\end{equation*}

Его изменение значительно слабее изменения множителя
$1/\beta^2$. Поэтому для качественной оценки можно считать

\begin{equation*}
    \rboxed{
        -
        \frac{dE}{d\ell}
        =
        \frac{C_1}{\beta^2}
    }.
\end{equation*}

По мере торможения протона его движение становится
нерелятивистским. Тогда кинетическая энергия равна

\begin{equation*}
    E
    =
    \frac{mv^2}{2}.
\end{equation*}

Так как

\begin{equation*}
    v=\beta c,
\end{equation*}

то

\begin{equation*}
    E
    =
    \frac{m\beta^2c^2}{2}.
\end{equation*}

Следовательно,

\begin{equation*}
    \beta^2
    =
    \frac{2E}{mc^2}.
\end{equation*}

Подставляя это выражение в

\begin{equation*}
    -
    \frac{dE}{d\ell}
    =
    \frac{C_1}{\beta^2},
\end{equation*}

получаем

\begin{equation*}
    -
    \frac{dE}{d\ell}
    =
    C_1
    \frac{mc^2}{2E}.
\end{equation*}

Все постоянные множители объединим в $C_2$:

\begin{equation*}
    \rboxed{
        -
        \frac{dE}{d\ell}
        =
        \frac{C_2}{E}
    }.
\end{equation*}

Отсюда

\begin{equation*}
    -E\,dE
    =
    C_2\,d\ell.
\end{equation*}

Интегрируем от начальной энергии $E_0$
до текущей энергии $E$:

\begin{equation*}
    -
    \int_{E_0}^{E}
    E\,dE
    =
    \int_0^\ell
    C_2\,d\ell.
\end{equation*}

Получаем

\begin{equation*}
    -
    \frac{
        E^2-E_0^2
    }{2}
    =
    C_2\ell.
\end{equation*}

Следовательно,

\begin{equation*}
    E_0^2-E^2
    =
    2C_2\ell.
\end{equation*}

Введём обозначение

\begin{equation*}
    C_3
    =
    2C_2.
\end{equation*}

Тогда

\begin{equation*}
    E_0^2-E^2
    =
    C_3\ell.
\end{equation*}

Отсюда энергия протона после прохождения расстояния $\ell$ равна

\begin{equation*}
    \rboxed{
        E(\ell)
        =
        \sqrt{
            E_0^2-C_3\ell
        }
    }.
\end{equation*}

Дифференцируя,

\begin{equation*}
    \frac{dE}{d\ell}
    =
    -
    \frac{
        C_3
    }{
        2
        \sqrt{
            E_0^2-C_3\ell
        }
    }.
\end{equation*}

Следовательно,

\begin{equation*}
    \rboxed{
        -
        \frac{dE}{d\ell}
        =
        \frac{
            C_3
        }{
            2
            \sqrt{
                E_0^2-C_3\ell
            }
        }
    }.
\end{equation*}

Таким образом, по мере уменьшения энергии частицы
её удельные энергетические потери возрастают.


\begin{rbox}[left=18pt,right=18pt,top=14pt,bottom=10pt]{[]}
    \centering

    \begin{tikzpicture}[
        >=Latex,
        every node/.style={text=rtext}
    ]

        % -----------------------------------------------------
        % Оси
        % -----------------------------------------------------

        \draw[
            ->,
            draw=rtext,
            line width=0.9pt
        ]
        (0,0) -- (7.0,0)
        node[right] {$\ell$};

        \draw[
            ->,
            draw=rtext,
            line width=0.9pt
        ]
        (0,0) -- (0,4.5)
        node[above] {$-\dfrac{dE}{d\ell}$};


        % -----------------------------------------------------
        % Брэгговская кривая
        % -----------------------------------------------------

        \draw[
            draw=rc3,
            line width=1.35pt,
            smooth
        ]
        plot coordinates {
            (0.25,0.62)
            (0.80,0.65)
            (1.40,0.69)
            (2.00,0.75)
            (2.60,0.84)
            (3.20,0.96)
            (3.75,1.13)
            (4.20,1.37)
            (4.60,1.72)
            (4.95,2.22)
            (5.22,2.92)
            (5.40,3.67)
            (5.50,4.05)
            (5.62,3.15)
            (5.78,2.05)
            (5.98,1.15)
            (6.20,0.48)
        };


        % -----------------------------------------------------
        % Положение пика
        % -----------------------------------------------------

        \draw[
            densely dashed,
            draw=rtextmuted,
            line width=0.7pt
        ]
        (5.50,0) -- (5.50,4.05);

        \draw[
            ->,
            draw=rtext,
            line width=0.8pt
        ]
        (4.15,3.65) -- (5.35,3.98);

        \node[
            above,
            align=center
        ] at (4.05,3.62) {
            брэгговский\\
            пик
        };

    \end{tikzpicture}

    \captionsetup{
        font=footnotesize,
        labelfont=bf,
        skip=5pt
    }

    \captionof{figure}{
        Качественная зависимость удельных энергетических потерь
        тяжёлой заряженной частицы от пройденного расстояния.
        Непосредственно перед остановкой частицы возникает
        брэгговский пик.
    }
    \label{fig:bragg_peak}
\end{rbox}


Основная причина такого поведения следует из зависимости

\begin{equation*}
    -
    \frac{dE}{d\ell}
    \propto
    \frac{1}{\beta^2}.
\end{equation*}

При торможении частицы

\begin{equation*}
    \beta
    \downarrow
    \qquad
    \Longrightarrow
    \qquad
    -
    \frac{dE}{d\ell}
    \uparrow.
\end{equation*}

Поэтому наиболее интенсивное энерговыделение происходит
в конце пробега тяжёлой заряженной частицы.

\end{document}
